\documentclass[../DoAn.tex]{subfiles}
\begin{document}

Chương 3 đã trình bày quá trình phân tích yêu cầu và thiết kế hệ thống. Chương này mô tả cách sinh viên đã hiện thực các module chính: môi trường phát triển, xác thực, tải lên và phân tích CV, trình soạn thảo CV, xuất CV, và luồng xử lý RAG tìm kiếm JD phù hợp cho ứng viên.

\section{Môi trường phát triển}

Hệ thống được phát triển bằng Next.js 15 \cite{nextjs2024} (App Router) với TypeScript \cite{typescript2024} và React 19 \cite{react2024} cho giao diện. Cơ sở dữ liệu dùng Supabase \cite{supabase2024} (PostgreSQL + pgvector \cite{pgvector2024}), lưu trữ tệp trên Vercel Blob \cite{vercel2024}, mô hình AI qua OpenAI API (\texttt{gpt-4o-mini} \cite{openai2023gpt4} và \texttt{text-embedding-3-small} \cite{openai2024embedding}). Xử lý kéo thả dùng \texttt{@dnd-kit}; trích xuất PDF dùng \texttt{pdf-parse} kết hợp \texttt{pdfjs-dist}; xuất DOCX dùng thư viện \texttt{docx}.

\section{Xác thực và quản lý phiên}

\textbf{Module xác thực} hỗ trợ hai chế độ: đăng nhập bằng email/mật khẩu (mật khẩu băm bằng \texttt{bcryptjs}, phiên duy trì bằng cookie bảo mật \texttt{httpOnly}) và phiên khách không cần tài khoản. Middleware Next.js kiểm tra cookie phiên trên mọi yêu cầu, cho phép truy cập editor ngay cả khi chưa đăng nhập -- dữ liệu khách lưu tạm trong trình duyệt. Khi khách đăng ký tài khoản, hệ thống nâng cấp dữ liệu nháp thành CV thuộc tài khoản mới mà không mất nội dung đang soạn.

\section{Trình soạn thảo CV}
\label{sec:editor_impl}

Trình soạn thảo là module trung tâm của sản phẩm -- nơi ứng viên nhìn thấy kết quả phân tích CV và chỉnh sửa nội dung trước khi xuất. Phần này mô tả bố cục giao diện, cơ chế lưu tự động, và tính năng kéo thả sắp xếp thứ tự.

\subsection{Bố cục giao diện hai cột}

Giao diện chia đôi màn hình theo chiều dọc: cột trái là \textit{editor} chứa các form chỉnh sửa từng section; cột phải là \textit{preview} hiển thị CV theo định dạng A4 cập nhật theo thời gian thực. Hình~\ref{fig:editor_layout} mô tả bố cục và các vùng chức năng chính.

\begin{figure}[H]
\centering
\includegraphics[width=\textwidth]{Hinh_ve/editor_layout.png}
\caption{Giao diện trình soạn thảo CV: editor hai cột (trái) và preview A4 cập nhật thời gian thực (phải)}
\label{fig:editor_layout}
\end{figure}

Mỗi section trong cột trái có một \textit{drag handle} (ba chấm dọc) ở cạnh trái -- kéo handle này để thay đổi thứ tự các section trên CV. Thư viện \texttt{@dnd-kit} quản lý trạng thái kéo thả và cập nhật mảng \texttt{sections[]} trong React state; thứ tự mới được phản ánh ngay trên preview A4 bên phải và ghi vào database khi auto-save kích hoạt.

\subsection{Cơ chế auto-save và đồng bộ preview}
\label{sec:autosave}

Thay vì yêu cầu người dùng nhấn nút ``Lưu'', trình soạn thảo tự động lưu sau mỗi lần chỉnh sửa. Hình~\ref{fig:autosave_seq} mô tả luồng từ thao tác gõ phím đến khi dữ liệu được ghi vào Supabase.

\begin{figure}[H]
\centering
\begin{tikzpicture}[
  life/.style={draw, fill=white, minimum width=2.2cm, minimum height=0.6cm,
               rounded corners=3pt, font=\small\bfseries, align=center},
  msg/.style={font=\scriptsize, align=center, inner sep=2pt},
  note/.style={draw, fill=yellow!20, rounded corners=3pt, font=\scriptsize,
               inner sep=5pt, text width=2.4cm, align=center},
  x=3.2cm, y=1cm,
]

%--- Lifeline headers (columns 0..4) ---
\node[life, fill=blue!12]   at (0,0) {Người dùng};
\node[life, fill=green!12]  at (1,0) {Editor UI};
\node[life, fill=orange!12] at (2,0) {Debounce\\Timer};
\node[life, fill=red!12]    at (3,0) {API autosave};
\node[life, fill=purple!12] at (4,0) {Supabase};

%--- Vertical lifelines ---
\foreach \c in {0,1,2,3,4}
  \draw[gray!40, densely dashed] (\c,-0.35) -- (\c,-8.8);

%--- Row -1.2 : user types ---
\draw[->, thick] (0.15,-1.2) -- (0.85,-1.2)
  node[midway, above, msg] {gõ phím};

%--- Row -2.2 : optimistic update (self-arrow on UI) ---
\draw[->, thick, blue!60] (1.15,-2.1) arc[start angle=0, end angle=-300, radius=0.18]
  node[right, msg, blue!80, text width=2.6cm] {cập nhật React state\\(preview ngay)};

%--- Row -3.2 : reset timer ---
\draw[->, thick] (1.15,-3.2) -- (1.85,-3.2)
  node[midway, above, msg] {reset timer 2s};

%--- Row -4.2 : user types again ---
\draw[->, thick] (0.15,-4.2) -- (0.85,-4.2)
  node[midway, above, msg] {gõ tiếp};

%--- Row -5.0 : reset timer again ---
\draw[->, thick] (1.15,-5.0) -- (1.85,-5.0)
  node[midway, above, msg] {reset timer 2s};

%--- Debounce waiting note ---
\node[note] at (2,-5.9) {Chờ 2s\\không có\\thay đổi};

%--- Row -6.8 : timer fires → API ---
\draw[->, thick, dashed] (2.15,-6.8) -- (2.85,-6.8)
  node[midway, above, msg] {POST /autosave};

%--- Row -7.6 : API → DB ---
\draw[->, thick] (3.15,-7.6) -- (3.85,-7.6)
  node[midway, above, msg] {UPDATE cv\_workflow};

%--- Row -8.2 : DB ack ---
\draw[->, thick, dashed] (3.85,-8.2) -- (3.15,-8.2)
  node[midway, above, msg] {200 OK};

%--- Row -8.8 : show saved ---
\draw[->, thick, dashed] (2.85,-8.8) -- (1.15,-8.8)
  node[midway, above, msg] {``Đã lưu lúc HH:MM''};

\end{tikzpicture}
\caption{Sequence diagram luồng auto-save với debounce 2 giây}
\label{fig:autosave_seq}
\end{figure}

Cơ chế debounce đảm bảo hệ thống không gọi API sau mỗi lần gõ phím -- chỉ gọi khi người dùng dừng gõ trong 2 giây. Preview phía phải cập nhật ngay lập tức từ React state (không chờ API), tạo cảm giác phản hồi tức thì. Nếu người dùng đóng tab trước khi debounce kích hoạt, \texttt{beforeunload} event gọi \texttt{navigator.sendBeacon} để gửi lần lưu cuối mà không block tab đóng.

\section{Xuất CV}

Hệ thống hỗ trợ hai định dạng xuất. \textbf{Xuất PDF} dùng \texttt{window.print()} với stylesheet \texttt{@media print} chuyên biệt: ẩn cột editor, đặt khổ A4, kiểm soát page-break để tránh section bị cắt đôi giữa trang. \textbf{Xuất DOCX} dùng thư viện \texttt{docx} tạo cấu trúc Word từ JSON CV -- mỗi section trong JSON được ánh xạ thành một block \texttt{Paragraph}/\texttt{Table} tương ứng trong file DOCX; font và spacing được chuẩn hoá để mở được trực tiếp bằng Microsoft Word mà không cần chỉnh sửa thêm.

\section{Hiện thực Hybrid CV Parser}
\label{sec:hybrid_parser_impl}

Khi ứng viên tải lên CV, văn bản thô từ PDF/DOCX thường có lỗi: chữ cái rời nhau, email bị ngắt dòng, tiêu đề mục lẫn vào nội dung. Hệ thống giải quyết bằng \texttt{EnhancedTextProcessor} -- pipeline tiền xử lý năm bước (chuẩn hoá ký tự, sửa lỗi PDF, chuẩn hoá tiêu đề mục, phát hiện ranh giới mục, đánh giá chất lượng văn bản) -- trước khi đưa vào parser. Bước phát hiện mục dùng ưu tiên tiêu đề rõ ràng; nếu CV không có tiêu đề chuẩn thì chuyển sang phân loại đoạn văn theo từ khoá đặc trưng.

Sau tiền xử lý, Hybrid CV Parser chạy hai pha song song. \textit{Pha 1} dùng regex trích các trường có cấu trúc cố định: email, số điện thoại Việt Nam (\texttt{+84}/\texttt{0} + 9--10 chữ số), URL LinkedIn, và kỹ năng kỹ thuật từ danh sách 10 nhóm từ khoá. \textit{Pha 2} gửi phần kinh nghiệm làm việc -- vốn viết tự do, khó tách bằng regex -- tới \texttt{gpt-4o-mini} với prompt yêu cầu trả về dữ liệu JSON có cấu trúc. Hai pha hợp nhất kết quả thành một JSON thống nhất; hệ thống tính thêm \texttt{possibility\_score} (0--100) để phát hiện tệp không phải CV. Cách kết hợp này giúp giảm lượng văn bản gửi lên LLM so với cách gửi nguyên văn bản thô, đồng thời giữ độ ổn định cao cho các trường có khuôn mẫu rõ.

\section{Hiện thực pipeline RAG tìm kiếm JD phù hợp}
\label{sec:rag_impl}

Khi ứng viên yêu cầu tìm JD khớp yêu cầu, hệ thống chạy quy trình hai bước qua đường dẫn API \texttt{/api/jobs/search}. Logic được tổ chức thành một module điều phối, gọi sang hai module phụ trợ riêng cho TF-IDF và Embedding ở bước retrieval.

\subsection{Bước 1 -- Hybrid Retrieval}

Mục tiêu của bước này là lọc nhanh từ toàn bộ kho JD xuống còn 20 tin tiềm năng, ưu tiên không bỏ sót hơn là chính xác. Hệ thống lấy vector ngữ nghĩa của CV (\texttt{cv\_embedding}) đã lưu sẵn trong bảng \texttt{cv\_workflow} -- không cần gọi thêm API -- rồi truy vấn bảng \texttt{job\_postings} để lấy toàn bộ vector JD đã được tính từ lúc crawl. Với mỗi JD, hệ thống tính điểm theo công thức~(\ref{eq:hybrid}): $\alpha = 0{,}4$ cho TF-IDF (bắt từ khoá kỹ thuật chính xác) cộng $(1-\alpha) = 0{,}6$ cho độ tương đồng ngữ nghĩa giữa hai vector. Hai nguồn điểm bổ sung cho nhau: TF-IDF tốt khi CV và JD dùng cùng từ ngữ; embedding tốt khi cùng ý nhưng diễn đạt khác nhau. Kết quả là top-20 JD điểm cao nhất được chuyển sang bước 2.

Vì cả CV embedding lẫn JD embedding đều được cache sẵn trong database, toàn bộ bước retrieval chạy thuần SQL -- không phát sinh chi phí API và độ trễ thấp do tận dụng được chỉ mục pgvector \cite{pgvector2024}.

\subsection{Bước 2 -- Structural Scoring và LLM Batch Ranking}

Bước 2 nhận 20 JD từ bước 1 và xếp hạng lại để trả về top-5 chính xác hơn. Điểm cuối cùng kết hợp hai thành phần theo công thức~(\ref{eq:final}).

\textbf{Structural Score} tính từ dữ liệu đã có sẵn trong JSON của CV, không cần gọi LLM. Ba tiêu chí được tính như sau:

\begin{itemize}
  \item \textbf{Kỹ năng} ($r$, trọng số 0,40): tỉ lệ giao giữa tập kỹ năng CV và top-20 từ khoá tần suất cao của JD -- $r = |\text{kỹ năng CV} \cap \text{kỹ năng JD}| \;/\; |\text{kỹ năng JD}|$.
  \item \textbf{Kinh nghiệm} ($e$, trọng số 0,35): điểm nhị phân -- 1 nếu tổng số năm kinh nghiệm cộng dồn từ \texttt{experience[].endDate - startDate} $\geq$ số năm yêu cầu trích từ JD bằng regex, ngược lại 0.
  \item \textbf{Bằng cấp} ($d$, trọng số 0,25): điểm nhị phân -- 1 nếu bậc học cao nhất trong \texttt{education[].degree} $\geq$ bậc yêu cầu trong JD (thang 0--4: trung cấp, cao đẳng, đại học, thạc sĩ, tiến sĩ), ngược lại 0.
\end{itemize}

Ba tiêu chí này được tính thuần trong bộ nhớ trong chưa đến 1 ms, không phát sinh chi phí API.

\textbf{LLM Batch Ranking} giải quyết những gì Structural Score bỏ qua -- ngữ cảnh công việc, văn hoá công ty, sự phù hợp tổng thể mà khó đo bằng số. Thay vì gọi LLM 20 lần riêng lẻ, hệ thống rút gọn mỗi JD về khoảng 200 đơn vị văn bản (giữ tiêu đề, kỹ năng, mức kinh nghiệm) và gom tất cả vào \textbf{một prompt duy nhất} cùng bản tóm tắt CV ứng viên. Mô hình \texttt{gpt-4o-mini} trả về danh sách xếp hạng dưới dạng JSON:

\begin{verbatim}
[{"rank": 1, "jobId": "job-42", "score": 87,
  "reason": "CV có 5 năm backend, thành thạo Node.js đúng yêu cầu JD"}]
\end{verbatim}

Cách gom nhóm này giảm số lần gọi API từ 20 xuống 1, tiết kiệm đáng kể chi phí so với gọi riêng lẻ từng JD.

\textbf{Tổng hợp và trả kết quả.} Điểm cuối cùng tính theo công thức~(\ref{eq:final}): $0{,}3 \cdot s_{\text{struct}} + 0{,}7 \cdot s_{\text{llm}}$. LLM được trọng số cao hơn vì nắm bắt được những yếu tố mềm mà Structural Score bỏ qua; Structural Score đóng vai trò neo, tránh LLM xếp hạng cao những JD quá xa tầm với của ứng viên. Top-5 JD kèm điểm thành phần và lý do giải thích được lưu vào \texttt{cv\_job\_searches.top\_k\_results} ở định dạng JSON để cache -- cùng một CV truy vấn lại không phải tính toán lại từ đầu.

\end{document}
